\label{sec:compactness-on-r}

\begin{topicbox}{The Open-Cover Definition}
\begin{exposition}
Compactness is defined by the open-cover property, deliberately and not as ``closed and bounded.'' The open-cover formulation is the one that survives to general metric spaces and to manifolds such as the torus, where ``bounded'' has no intrinsic meaning. On $\mathbb{R}$ the two descriptions coincide, and that coincidence is the Heine--Borel theorem of the next section, proved rather than assumed.
\end{exposition}
\end{topicbox}
\begin{definitionbox}{Definition (Compact Set in $\mathbb{R}$)}
\begin{definition}[Compact Set in $\mathbb{R}$]
\label{def:compact-set}
A set $K\subseteq\mathbb{R}$ is \emph{compact} if every open cover of $K$ has
a finite subcover.  Explicitly, for every family $\mathcal{U}$ of open
subsets of $\mathbb{R}$,
\[
  K\subseteq\bigcup_{U\in\mathcal{U}}U
\]
implies that there exists a finite subfamily $\mathcal{V}\subseteq\mathcal{U}$
such that
\[
  K\subseteq\bigcup_{V\in\mathcal{V}}V.
\]
\end{definition}
\end{definitionbox}
\begin{remark*}[Standard quantified statement]
\[
K \text{ compact} \Longleftrightarrow \forall \{U_i\}_{i\in I}\Bigl(\bigl(\forall i\in I\;U_i \text{ open}\bigr)\land K\subseteq\textstyle\bigcup_{i\in I}U_i \Longrightarrow \exists \text{finite } J\subseteq I\;K\subseteq\textstyle\bigcup_{i\in J}U_i\Bigr).
\]
\end{remark*}

\begin{remark*}[Predicate reading]
\[
K\subseteq\mathbb{R},\qquad
\operatorname{IsCompact}(K,\mathbb{R})
\Longleftrightarrow
\forall \mathcal{U}\;\bigl(\operatorname{OpenCover}(\mathcal{U},K,\mathbb{R})
\Longrightarrow \exists J\;(J\subseteq I\land \operatorname{FiniteSet}(J)
\land K\subseteq\textstyle\bigcup_{i\in J}U_i)\bigr).
\]
\end{remark*}
\begin{remark*}[Negated quantified statement]
\[
K \text{ not compact} \Longleftrightarrow \exists \{U_i\}_{i\in I}\Bigl(\bigl(\forall i\in I\;U_i \text{ open}\bigr)\land K\subseteq\textstyle\bigcup_{i\in I}U_i \land \forall \text{finite } J\subseteq I\;K\not\subseteq\textstyle\bigcup_{i\in J}U_i\Bigr).
\]
\end{remark*}
\begin{remark*}[Negation predicate reading]
\[
K\subseteq\mathbb{R},\qquad
\neg\operatorname{IsCompact}(K,\mathbb{R})
\Longleftrightarrow
\exists \mathcal{U}\;\bigl(\operatorname{OpenCover}(\mathcal{U},K,\mathbb{R})
\land \forall J\;(J\subseteq I\land \operatorname{FiniteSet}(J)
\Longrightarrow K\not\subseteq\textstyle\bigcup_{i\in J}U_i)\bigr).
\]
\end{remark*}

\begin{remark*}[Failure modes]
\begin{description}
  \item[Exposition.]
  Failure is witnessed by the displayed negated or contrapositive condition.

  \item[\operatorname{IsCompact} / \operatorname{OpenCover}.]
  \[
K \text{ not compact} \Longleftrightarrow \exists \{U_i\}_{i\in I}\Bigl(\bigl(\forall i\in I\;U_i \text{ open}\bigr)\land K\subseteq\textstyle\bigcup_{i\in I}U_i \land \forall \text{finite } J\subseteq I\;K\not\subseteq\textstyle\bigcup_{i\in J}U_i\Bigr).
\]
\[
K\subseteq\mathbb{R},\qquad
\neg\operatorname{IsCompact}(K,\mathbb{R})
\Longleftrightarrow
\exists \mathcal{U}\;\bigl(\operatorname{OpenCover}(\mathcal{U},K,\mathbb{R})
\land \forall J\;(J\subseteq I\land \operatorname{FiniteSet}(J)
\Longrightarrow K\not\subseteq\textstyle\bigcup_{i\in J}U_i)\bigr).
\]
\end{description}
\end{remark*}
\begin{remark*}[Interpretation]
Compactness is a finiteness property phrased entirely in open sets: no matter how a set is covered by open pieces, finitely many already do the job. This is what later guarantees, on a compact domain, that continuous functions attain extrema and are uniformly continuous --- the facts Analysis~I relies on. Crucially, this definition does not mention order or boundedness, so it transfers unchanged to the torus.
\end{remark*}
\begin{dependencies}
\begin{itemize}
  \item \hyperref[def:real-open-cover]{Open Cover}
  \item \hyperref[def:finite-subcover]{Subcover and Finite Subcover}
\end{itemize}
\end{dependencies}
\begin{theorembox}{Theorem (Compact Subsets of $\mathbb{R}$ Are Closed and Bounded)}
\begin{theorem}[Compact Subsets of $\mathbb{R}$ Are Closed and Bounded]
\label{thm:compact-implies-closed-bounded}
Let $K\subseteq\mathbb{R}$.  If $K$ is compact, then $K$ is closed and
bounded.  Explicitly,
\[
  K\text{ compact}
  \Longrightarrow
  \leqft[
  \mathbb{R}\setminus K\text{ is open}
  \right]
  \land
  \leqft[
  (\exists M>0)(\forall x\in K)(|x|\leq M)
  \right].
\]
\hyperref[prf:compact-implies-closed-bounded]{Go to proof.}
\end{theorem}
\end{theorembox}
\begin{remark*}[Standard quantified statement]
\[
K \text{ compact} \Longrightarrow (K \text{ closed} \land K \text{ bounded}).
\]
\end{remark*}
\begin{remark*}[Interpretation]
This is the easy half of Heine--Borel and holds the same way in every metric space. Boundedness comes from covering $K$ by growing balls about a fixed point; closedness comes from separating an outside point from $K$ by disjoint balls. The converse --- closed and bounded implies compact --- is special to $\mathbb{R}^{n}$ and is the next theorem.
\end{remark*}
\begin{dependencies}
\begin{itemize}
  \item \hyperref[def:compact-set]{Compact Set in $\mathbb{R}$}
  \item \hyperref[def:closed-set]{Closed Set in $\mathbb{R}$}
  \item \hyperref[def:bounded-subset-of-r]{Bounded Subset of the Real Line}
\end{itemize}
\end{dependencies}
